\begin{pastexam}[label={px:nccu-108-dsa-08}]{政大資科所 108 學年度 第 8 題}
（問答）

(10\%) Give a binary heap containing the following data: 1, 2, 3, 4, 5, 6, 7.
\end{pastexam}

\begin{pastsoln}
\textbf{解讀}。 題目沒說最小堆積還是最大堆積，答案也不唯一，給出一個合法的堆積並檢查兩個條件即可：
(1) 完全二元樹；(2) 堆積性質（資結書第~8.1~節）。 下面兩種都寫。

\textbf{最小堆積}。 由小到大排好的陣列本身就是最小堆積：$A = \langle 1, 2, 3, 4, 5, 6, 7 \rangle$，
節點 $i$ 的孩子是 $2i$、$2i + 1$，孩子的索引比較大，值也比較大。
\begin{center}
\begin{tikzpicture}[x=1cm,y=1cm,v/.style={circle,draw,minimum size=6mm,inner sep=0pt,font=\small}]
\begin{scope}
\node[v] (a) at (0,0) {1};
\node[v] (b) at (-1.5,-1) {2}; \node[v] (c) at (1.5,-1) {3};
\node[v] (d) at (-2.2,-2) {4}; \node[v] (e) at (-0.8,-2) {5}; \node[v] (f) at (0.8,-2) {6}; \node[v] (g) at (2.2,-2) {7};
\draw (a)--(b) (a)--(c) (b)--(d) (b)--(e) (c)--(f) (c)--(g);
\node[font=\small] at (0,-2.8) {最小堆積 $\langle 1, 2, 3, 4, 5, 6, 7 \rangle$};
\end{scope}
\begin{scope}[xshift=6.5cm]
\node[v] (a) at (0,0) {7};
\node[v] (b) at (-1.5,-1) {5}; \node[v] (c) at (1.5,-1) {6};
\node[v] (d) at (-2.2,-2) {4}; \node[v] (e) at (-0.8,-2) {2}; \node[v] (f) at (0.8,-2) {1}; \node[v] (g) at (2.2,-2) {3};
\draw (a)--(b) (a)--(c) (b)--(d) (b)--(e) (c)--(f) (c)--(g);
\node[font=\small] at (0,-2.8) {最大堆積 $\langle 7, 5, 6, 4, 2, 1, 3 \rangle$};
\end{scope}
\end{tikzpicture}
\end{center}
檢查（最小堆積）：$1 \le 2, 3$；$2 \le 4, 5$；$3 \le 6, 7$。 $7$ 個節點剛好填滿三層，是完全二元樹。

\textbf{最大堆積}。 把 $\langle 1, 2, 3, 4, 5, 6, 7 \rangle$ 用由下而上建堆積（資結書虛擬碼~8.3，
\proc{Build-Max-Heap}）：$i$ 從 $\lfloor 7/2 \rfloor = 3$ 往下到 $1$，對每個 $i$ 下沉。
\begin{center}
\begin{tabular}{cl}
\toprule
$i$ & 下沉後的陣列 \\
\midrule
3 & $\langle 1, 2, \mathbf{7}, 4, 5, 6, \mathbf{3} \rangle$（$3$ 和較大的孩子 $7$ 交換） \\
2 & $\langle 1, \mathbf{5}, 7, 4, \mathbf{2}, 6, 3 \rangle$（$2$ 和較大的孩子 $5$ 交換） \\
1 & $\langle \mathbf{7}, 5, \mathbf{1}, 4, 2, 6, 3 \rangle \to \langle 7, 5, \mathbf{6}, 4, 2, \mathbf{1}, 3 \rangle$（$1$ 下沉兩層） \\
\bottomrule
\end{tabular}
\end{center}
得到 $\langle 7, 5, 6, 4, 2, 1, 3 \rangle$（右上圖）。 檢查：$7 \ge 5, 6$；$5 \ge 4, 2$；$6 \ge 1, 3$。
最簡單的最大堆積是由大到小的 $\langle 7, 6, 5, 4, 3, 2, 1 \rangle$，也對。

注意堆積只要求父節點不小於（或不大於）孩子，不要求左孩子小於右孩子，這和二元搜尋樹不同；
同一份考卷第 7 題（歷屆試題~\ref{px:nccu-108-dsa-07}）的二元搜尋樹 $\langle 4, 2, 6, 1, 3, 5, 7 \rangle$（依層序）就不是堆積（$4 > 2$ 但 $4 < 6$）。
\end{pastsoln}
